\documentclass[dvipdfmx]{jsarticle}
\usepackage[margin=1cm]{geometry}
\usepackage{amsmath}

\begin{document}

\fboxsep0pt
\fbox{%
\begin{minipage}[t][17.5cm][t]{17.5cm}
\vspace{0.5cm}
\hspace{0.5cm}%
\begin{minipage}{16.5cm}


\noindent
\textbf{1.\ 北海道大学理学部数学科への進学を志望する背景}

 
\smallskip
\noindent
北海道大学理学院数学専攻を志望するのは，数学の中でも様々な分野の教員・学生が在籍しているからです．私の志望する数理科学系が盛んであるのも珍しいですが，かといって代数や幾何の教員が少ないというわけではないため，他分野の方と交流することで様々なトピックに触れることができ，自分の分野の理解も広げてくれるだろうと思います．とくに北大は若手研究集会やシンポジウムなどもやっていますし，よく任意参加のセミナーも開催されているので，そのような横断的な交流がしやすいと思っています．
 
\bigskip
\noindent
\textbf{2.\ 希望する系（数理科学系）とその理由}
 
\smallskip
\noindent
私が希望するのは数理科学系になります．
代数・幾何が主に数学的対象そのものの構造を追究するのに対し，数理科学系は確率論・統計学・数理物理学など，具体的な現象と結びついた数学を扱う点に特色があると理解しています．
学部3年までどの分野も興味深く学んでいましたが，3年にルベーグ積分論（測度論）と確率論の講義を受け，確率を測るという発想やルベーグの収束定理の証明などに感銘を受け，この分野に進もうと思うようになりました．また統計学の講義ではその発想，および実用性に魅力を感じ，4年の今は工学部のデータ分析の講義を取って実社会での応用を興味深く学んでいます．このような，ある程度現実の事象との結びつきを感じられる数学を学びたいと思い，数理科学系を志望しました．
また，これまで数学的な基礎を固めたことですし，完全な応用に傾くというよりかは厳密にロジックを築いて自分の興味のある分野の全体像の解像度をあげていきたいと思っています．そのため，工学院ではなく理学院を希望します．
 
\bigskip
\noindent
\textbf{3.\ 具体的に研究したいこと}
 
\smallskip
\noindent
統計物理学・場の量子論に由来する「くりこみ群」的アプローチが，確率論における極限定理とどのように結びつくのかに強い関心があります．Koralov-SinaiやOttによるくりこみ群と中心極限定理の関係を論じた文献を手がかりに，スケール変換のもとでの普遍性という物理的な概念を，確率論的な極限定理の枠組みの中で厳密に理解していきたいと考えています．
また，その後はくりこみ群の本来の主戦場である統計物理学の格子模型，とりわけイジングモデルやパーコレーション，自己回避歩行といった系の臨界現象へと視野を広げたいと考えています．

\end{minipage}
\end{minipage}%
}

\end{document}